\section{Design Considerations}
\label{sec:hw_design}

The hardware design of the Vayu flight control system is guided by several key considerations that ensure reliable operation, efficient performance, and seamless integration with the software architecture.

\subsection{Real-Time Constraints}

Flight control systems operate under strict real-time requirements. The hardware is selected to support deterministic execution through the availability of high-resolution timers, low-latency interrupt handling, and efficient communication interfaces. These features enable consistent sampling of sensor data and timely execution of control loops, which are essential for system stability.

Additionally, the allocation of dedicated interfaces for time-critical operations—such as UART for radio input and timers for actuation—ensures minimal latency and predictable system behavior.

\subsection{Signal Integrity and Noise Management}

Maintaining signal integrity is critical, particularly for sensitive components such as the IMU. The use of a two-stage power regulation system (buck converter followed by LDO), along with proper decoupling and grounding practices, helps reduce electrical noise. This ensures that sensor measurements remain accurate and are not adversely affected by power supply fluctuations or electromagnetic interference.

Special attention is given to PCB layout, component placement, and power routing to minimize noise coupling between high-current switching elements and precision sensing circuits.

\subsection{Modularity and Expandability}

The hardware platform is designed with modularity in mind, exposing communication interfaces such as UART, SPI, and I2C through dedicated connectors. This allows additional sensors, peripherals, and communication modules to be integrated without redesigning the core system.

Future revisions will extend this modularity further through additional interfaces such as USB-based communication and expanded bus availability, enabling integration with companion systems and advanced modules.

\subsection{Scalability and Upgrade Path}

The system is designed to accommodate future enhancements in both hardware and software. The use of a hardware abstraction layer (NavHAL) enables portability across different microcontrollers, while the modular hardware design supports the addition of redundant sensors, advanced communication protocols, and higher-performance processing units.

Planned upgrades include higher-performance MCUs, SPI-based high-speed sensors, and digital ESC communication protocols such as DShot.

\subsection{Debugging and Development Support}

Development and testing are facilitated through dedicated debugging interfaces such as SWD and UART-based telemetry. These interfaces allow real-time monitoring of system behavior, aiding in debugging, performance tuning, and validation of control algorithms.

Future iterations will incorporate USB CDC for higher-bandwidth debugging and interaction with external computing systems.

\subsection{Hardware–Software Alignment}

A key consideration throughout the design is the alignment between hardware capabilities and software requirements. The selection of peripherals, communication interfaces, and timing resources is closely matched to the needs of NavHAL and VAIOS, ensuring efficient abstraction and deterministic execution.

\subsection{System Reliability and Safety}

The hardware design is structured to support reliable system operation under dynamic conditions. Stable power regulation, proper grounding, and isolation of critical interfaces ensure consistent behavior across varying load scenarios.

Future updates will incorporate battery monitoring capabilities, including voltage and current sensing, enabling implementation of power-aware safety mechanisms such as low-battery failsafe and controlled landing strategies.

Overall, these design considerations ensure that the hardware platform not only supports current system requirements but also provides a robust, scalable, and reliable foundation for future development and deployment of the Vayu flight control stack.